% Part: incompleteness
% Chapter: incompleteness-provability
% Section: lob-thm

\documentclass[../../../include/open-logic-section]{subfiles}

\begin{document}

\olfileid{inc}{inp}{tar}

\olsection{عدم قابلية الصدق للتعريف}

يعتمد مفهوم \emph{قابلية التعريف} على وجود دلالة صورية للغة الحساب.
وقد وصفنا مجموعة من الصيغ والجمل في لغة الحساب. ويتمثل «التأويل
المقصود» في قراءة هذه الجمل على أنها تقرر أمورًا عن الأعداد الطبيعية،
ويمكن لمثل هذا التقرير أن يكون صادقًا أو كاذبًا. ولتكن~$\Struct{N}$
هي الـ!!{structure} ذات المجال~$\Nat$ والتأويل القياسي لرموز لغة
الحساب. عندئذ تعني~$\Sat{N}{!A}$ أن «$!A$ صادقة في التأويل القياسي».

\begin{defn}
تكون العلاقة~$R(x_1,\dots,x_k)$ بين الأعداد الطبيعية
\emph{قابلة للتعريف} في~$\Struct{N}$ إذا وفقط إذا وجدت صيغة
$!A(x_1,\dots,x_k)$ في لغة الحساب بحيث، لكل
$n_1,\dots,n_k$، تصدق~$R(n_1,\dots,n_k)$ إذا وفقط إذا
$\Sat{N}{!A(\num n_1,\dots,\num n_k)}$.
\end{defn}

وبعبارة أخرى، تكون العلاقة قابلة للتعريف في~$\Struct{N}$ إذا وفقط
إذا كانت قابلة للتمثيل في النظرية~$\Th{TA}$، حيث
$\Th{TA} = \Setabs{!A}{\Sat{N}{!A}}$ هي مجموعة جمل الحساب الصادقة.
(وإذا لم يكن هذا واضحًا لك فورًا، فينبغي أن تعود إلى التعريفات وتتحقق
بنفسك من صحة ذلك.)

\begin{lem}
كل علاقة قابلة للحساب قابلة للتعريف في~$\Struct{N}$.
\end{lem}

\begin{proof}
يسهل التحقق من أن الصيغة التي تمثل علاقة في~$\Th{Q}$ تعرّف العلاقة
نفسها في~$\Struct{N}$.
\end{proof}

يمكننا الآن أن نسأل: هل يصدق العكس أيضًا؟ أي هل كل علاقة قابلة
للتعريف في~$\Struct{N}$ قابلة للحساب؟ الجواب لا. ومن أمثلة ذلك:

\begin{lem}
علاقة التوقف قابلة للتعريف في~$\Struct{N}$.
\end{lem}

\begin{proof}
تذكّر أن مبرهنة الصورة السوية لكلييني تقرر أن لكل دالة جزئية قابلة
للحساب~$f$ دليلًا~$e$ بحيث
$f(x) = U(\umin{s}{T(e,x,s)})$ لكل~$x \in \Nat$، حيث الدالتان
$U$ و$T$ عوديتان بدائيتان، ومن ثم كليتان. ولذلك تكون~$f(x)$ معرفة
(أي تتوقف العملية الحسابية) إذا وفقط إذا وجد~$s$ بحيث تصدق
$T(e,x,s)$.

لتكن الآن~$H$ علاقة التوقف؛ أي
\[
H = \Setabs{\tuple{e,x}}{\lexists[s][T(e, x, s)]}.
\]
ولتكن~$!D_T$ تعرّف~$T$ في~$\Struct{N}$. عندئذ
\[
H = \Setabs{\tuple{e,x}}{\Sat{N}{\lexists[s][!D_T(\num e, \num x, s)]}},
\]
ومن ثم تعرّف~$\lexists[s][!D_T(z, x, s)]$ العلاقة~$H$ في
$\Struct{N}$.
\end{proof}

\begin{prob}
بيّن أن~$Q(n) \defiff n \in \Setabs{\Gn{!A}}{\Th{Q} \Proves !A}$
قابلة للتعريف في الحساب.
\end{prob}

وماذا عن~$\Th{TA}$ نفسها؟ هل هي قابلة للتعريف في الحساب؟ أي هل
المجموعة~$\Setabs{\Gn{!A}}{\Sat{N}{!A}}$ قابلة للتعريف في الحساب؟
تجيب مبرهنة تارسكي بالنفي.

\begin{thm}
\ollabel{thm:tarski}
مجموعة !!{sentence}s الحسابية الصادقة غير قابلة للتعريف في الحساب.
\end{thm}

\begin{proof}
افترض أن~$!D(x)$ تعرّفها؛ أي إن~$\Sat{N}{!A}$ إذا وفقط إذا
$\Sat{N}{!D(\gn{!A})}$. وبحسب لمّة النقطة الثابتة، توجد جملة~$!A$
بحيث~$\Th{Q} \Proves !A \liff \lnot !D(\gn{!A})$، ومن ثم
$\Sat{N}{!A \liff \lnot !D(\gn{!A})}$. ولكن عندئذ تكون
$\Sat{N}{!A}$ إذا وفقط إذا~$\Sat{N}{\lnot !D(\gn{!A})}$، وهذا يناقض
كون~$!D(y)$ مفترضةً أن تعرّف مجموعة عبارات الحساب الصادقة.
\end{proof}

طبق تارسكي هذا التحليل على مفهوم فلسفي أعم للصدق. فبالنسبة إلى أي
لغة~$L$، رأى تارسكي أن مفهومًا ملائمًا للصدق في~$L$ يجب أن يحقق، لكل
جملة~$X$، ما يأتي:
\begin{quote}
«$X$» صادقة إذا وفقط إذا~$X$.
\end{quote}
ومثال تارسكي كثير الاقتباس، في الإنجليزية، هو الجملة:
\begin{quote}
«الثلج أبيض» صادقة إذا وفقط إذا كان الثلج أبيض.
\end{quote}
غير أننا نستطيع، لكل لغة قوية بما يكفي لتمثيل الدالة القطرية ولكل
محمول لغوي~$T(x)$، أن نبني جملة~$X$ تحقق: «$X$ إذا وفقط إذا لم تكن
$T(\text{«$X$»})$». ولما كنا لا نريد لمحمول الصدق أن يعلن أن بعض
الجمل صادقة وكاذبة معًا، خلص تارسكي إلى أنه لا يمكن تحديد محمول صدق
لجميع جمل لغة ما من دون الخروج، على نحو ما، عن حدود تلك اللغة. أي إن
محمول الصدق للغة لا يمكن تعريفه في اللغة نفسها.

\end{document}
